\chapter{Project Proposal}
\ChapterLabel{ProjectProposal}
\index{project proposal}
\index{DevSecOps}
\index{DevSecOps!tool candidates}

\section{Project Description}

\subsection{Working Title}
\textbf{StudyOps: A DevSecOps-Built Personal Academic Dashboard.}

\subsection{Description}
StudyOps is a small, self-hosted web application that gives a single place to
track every course, assignment, exam, and study session for the semester,
replacing a scattered mix of course syllabi, spreadsheets, and sticky notes.
Functionally it is a three-tier application: a lightweight browser frontend,
a REST API backend, and a relational database holding each course,
assignment, and study-session record for the logged-in student. The
application itself is intentionally simple -- the point of the project is
the DevSecOps pipeline and environment built around it, not a complicated
feature set -- but it is also something I would genuinely use for my own
semester planning, which keeps me motivated to keep it running correctly
rather than letting it become shelfware once the assignment is graded.

The backend will be written in \textbf{C\# with ASP.NET Core}, a deliberate
choice: it is one of the most common backend stacks in corporate/enterprise
settings, and neither of us has hands-on background in it, so the project
doubles as an opportunity to build that skill before graduation. The
frontend is kept intentionally simple (plain HTML/JavaScript calling the
API) so the new-to-us learning stays focused on the backend, data access,
and DevSecOps pipeline rather than also picking up an unfamiliar frontend
framework at the same time.

Everything below is a first pass. The specific tool names are candidates we
intend to evaluate, not final commitments; several categories list more than
one option because part of the project is comparing them directly (see
Section~\ref{Section:ToolComparison}).

\subsection{Sample Tasks}
At minimum, the finished application should support:
\begin{itemize}
\item Creating a course (name, term, credit hours).
\item Adding an assignment or exam to a course with a due date and a priority/weight.
\item Viewing a single ``what's due'' dashboard across all courses, sorted soonest-first.
\item Marking an assignment complete, and logging a study session against a course.
\item Receiving a generated weekly digest of everything due in the next 7 days (described further in Section~\ref{Section:Unique}).
\end{itemize}

\subsection{Candidate Toolchain}
\begin{tabularx}{\linewidth}{L{3.0cm}L{3.6cm}Y}
\caption{Candidate DevSecOps Toolchain\label{Table:ToolchainCandidates}}\\
\toprule \textbf{Category} & \textbf{Candidate tool(s)} & \textbf{Notes}\\ \midrule
Application stack & Backend: C\# / ASP.NET Core Web API. Frontend: plain HTML/JavaScript (kept simple on purpose) & Chosen for industry relevance in corporate/enterprise backends; neither of us has prior C\#/.NET background, so this is a deliberate skill-building choice.\\
Data access & Entity Framework Core (ORM) against PostgreSQL & Keeps schema changes (migrations) version-controlled alongside the C\# code.\\
Source control & GitHub & Private repository, branch protection on \texttt{main}, pull-request review required before merge.\\
CI/CD pipeline & GitHub Actions; compared against GitLab CI & Build, test, scan, and deploy on every push with the \texttt{dotnet} CLI; the second tool is the ``comparable tool'' pairing for the pipeline layer.\\
Containerization & Docker, built from the official Microsoft ASP.NET Core runtime image, with Docker Compose for local multi-container development (API + database) & Same image is promoted from local dev through the pipeline to each hosting environment.\\
Container orchestration (stretch) & Lightweight Kubernetes (k3s), compared against plain Docker Compose on a single host & Only attempted if time allows after the core pipeline is working.\\
Infrastructure as Code & Terraform; compared against Pulumi if time allows & Provisions the compute, network, and database resources in each cloud environment from version-controlled code.\\
Hosting environment & AWS (EC2/ECS + RDS) versus DigitalOcean (Droplets/App Platform + Managed Database) & The project's primary two-environment comparison; same container image deployed to both. Both platforms run containerized .NET apps without issue.\\
Database & PostgreSQL, run as a managed service on each platform (RDS vs.\ DigitalOcean Managed Databases) & Keeps the comparison consistent through the data layer.\\
Automated testing & xUnit for unit tests, plus a small integration-test suite run with \texttt{dotnet test} against the live API in the pipeline & Tests gate the deploy step; a failing test blocks release.\\
Security scanning & Dependabot or Snyk (dependency/NuGet-package vulnerabilities), Trivy (container-image scanning), gitleaks or truffleHog (committed-secret scanning) & Run as pipeline gates, the ``Sec'' in DevSecOps for this project.\\
Secrets management & Pipeline-native encrypted secrets (GitHub Actions secrets) and per-environment \texttt{.env}/\texttt{appsettings} files, never committed & HashiCorp Vault considered as a stretch goal.\\
Monitoring / observability & CloudWatch (AWS) versus DigitalOcean Monitoring (DO), or a self-hosted Prometheus + Grafana stack for a platform-neutral view & Feeds the operational dashboard in Section~\ref{Section:Dashboard}.\\ \bottomrule
\end{tabularx}

\section{Planned Tool Comparison}
\label{Section:ToolComparison}
The project's primary ``two comparable tools'' comparison is at the hosting-environment
layer: the identical container image is deployed to both \textbf{AWS} and
\textbf{DigitalOcean}, and the two are compared directly across:
\begin{itemize}
\item \textbf{Security} -- identity and access model (IAM vs.\ DigitalOcean's simpler access model), network/firewall configuration, and how cleanly each platform keeps secrets out of version control.
\item \textbf{Development experience} -- CLI/console ergonomics, documentation quality, and time to a first working deploy.
\item \textbf{Hosting} -- pricing model, free-tier limits, and managed-database setup effort.
\item \textbf{Monitoring} -- built-in dashboards, alerting, and log retention.
\item \textbf{Testing} -- how (and where) the pipeline's test and security-scan stages differ, if at all, per target environment.
\item \textbf{Operations} -- deploy frequency achievable in practice, rollback process, and general day-2 operations burden.
\end{itemize}
If time allows, a secondary comparison at the CI/CD layer (GitHub Actions vs.\
GitLab CI) will be attempted using the same application and the same AWS
target, to isolate the pipeline tool's effect from the hosting environment's.

\section{Environment}
The project will be built using a real cloud environment rather than a
purely local (``homebrew'') setup, specifically because comparing AWS and
DigitalOcean is the centerpiece of the comparison plan above. Both platforms
offer a free or low-cost tier sufficient for a single small application, and
local Docker Compose remains the development and first-test environment
before anything is deployed.

\section{Operational Dashboard (Planned)}
\label{Section:Dashboard}
One operational dashboard per hosting environment (or a single unified
dashboard if a self-hosted Prometheus/Grafana stack is used instead) should
reflect, at minimum:
\begin{itemize}
\item Deployment frequency and lead time (time from commit to live).
\item Build/test/scan pass-fail history for recent pipeline runs.
\item Current application health (uptime / health-check status).
\item Error rate and response latency for the API.
\item The commit SHA and trigger (push, scheduled job, manual) of the most recent deploy.
\end{itemize}

\section{Preliminary SWOT (First Pass)}
These SWOT charts are an early, intentionally rough first pass on the two
hosting candidates, to be refined once hands-on work begins; they are
included now to show the comparison is already being thought through rather
than deferred.

\subsection{AWS}
\begin{tabularx}{\linewidth}{L{3.2cm}Y}
\toprule \textbf{Category} & \textbf{Notes}\\ \midrule
Strengths & Broadest service catalog and the cloud platform most referenced in job postings; mature IAM and security tooling; a generous, well-documented free tier for small EC2/RDS instances.\\
Weaknesses & Steeper learning curve; pricing and service sprawl make cost estimation harder for a single small app; the console is dense for a first-time user.\\
Opportunities & Direct experience with the most commonly requested cloud skill in corporate job postings; room to later add managed CI (CodePipeline) or a serverless function (Lambda) for the weekly-digest feature.\\
Threats & Easy to accidentally leave a billable resource running past the free tier; a misconfigured IAM policy or security group is a realistic risk on a first AWS project.\\ \bottomrule
\end{tabularx}

\subsection{DigitalOcean}
\begin{tabularx}{\linewidth}{L{3.2cm}Y}
\toprule \textbf{Category} & \textbf{Notes}\\ \midrule
Strengths & Simple, flat pricing and a small, approachable product set (Droplets, App Platform, Managed Databases); fast path to a first working deploy; documentation aimed at smaller projects.\\
Weaknesses & Smaller service catalog, so some AWS comparisons (fine-grained IAM, serverless) have no direct equivalent; less name recognition on a resume.\\
Opportunities & A faster path to a working deployment frees up more project time for the DevSecOps pipeline itself; App Platform's built-in build-from-Git-push is a natural comparison point against a hand-built GitHub Actions pipeline.\\
Threats & Smaller platform means fewer community answers for an obscure error; less corporate adoption than AWS, so the experience is somewhat less directly transferable to a specific employer's stack.\\ \bottomrule
\end{tabularx}

\section{Something Unique}
\label{Section:Unique}
\subsection{Weekly Deadline Digest}
A scheduled pipeline job (a GitHub Actions cron trigger, or a small
serverless function on a timer) queries StudyOps's own database once a week
and sends a short digest -- everything due in the next 7 days -- to a
personal email address or a Discord/Slack webhook. This turns the
``deploy and monitor'' lesson into something with a direct personal payoff,
and it is also a deliberately different trigger model (scheduled,
event-driven automation) from the push-triggered pipeline used for deploys,
so the project demonstrates more than one automation pattern.

\subsection{Live Pipeline-Status Badge}
A second, smaller unique touch: a public, read-only ``last deploy:
passed/failed'' status badge in the project's README, the same convention
many open-source projects use, adapted here to a personal project's own
pipeline.

\section{Expectations and Rubric}
The project will be evaluated on:
\begin{itemize}
\item The selected tool, tools, or environment.
\item Building a sample working application.
\item Using specific DevOps tools to build a simple project.
\item Ideally, using two comparable tools to build the same project and drawing comparisons across security, development, hosting, monitoring, testing, and operations.
\item Using a real environment to build the project -- AWS, DigitalOcean, or a homebrew setup (containers, microservices, GitHub, etc.).
\item Constructing a dashboard reflecting operational aspects.
\item Constructing a modified Strengths/Weaknesses/Opportunities/Threats (SWOT) chart for each tool.
\item Proposing something unique.
\end{itemize}

\begin{devopsNote}[Status]
This chapter is a proposal. Tool names, the specific comparison pairs, and
the unique-feature ideas above are the group's current best thinking and are
expected to be refined (or changed outright) once implementation starts;
later chapters will record what was actually built, with evidence, the same
way \ChapterRef{LinuxCommands} does for the problem-set work. The backend
language/framework (C\#/ASP.NET Core) is a firmer commitment than the other
candidates, since it was chosen specifically to build skills neither of us
currently has; the hosting, CI/CD, and IaC tool pairings remain open.
\end{devopsNote}
